\newpage
\begin{center}
  \textbf{\large ПРИЛОЖЕНИЕ А}\\
  \textbf{\large Чек-лист предпроектного обследования клиентского цифрового контура}
\end{center}
\refstepcounter{chapter}
\addcontentsline{toc}{chapter}{ПРИЛОЖЕНИЕ А. Чек-лист предпроектного обследования клиентского цифрового контура}
\label{app:ch2_audit}

Настоящее приложение фиксирует первичный материал, на основании которого во второй главе рассчитаны показатели AS-IS состояния клиентского цифрового контура ООО <<\_\_\_\_>>. Чек-лист не заменяет бухгалтерскую, управленческую или маркетинговую отчетность предприятия, поскольку объект находится на раннем этапе развития и не имеет накопленной статистики цифровых заявок. Его назначение уже: подтвердить измеримые факты предпроектного обследования, связанные с наличием каналов, доступностью информации, структурированностью обращения и возможностью дальнейшего анализа клиентского опыта.

\begin{table}[H]
\centering
\small
\caption{Паспорт предпроектного обследования}
\label{tab:app_audit_passport}
\begin{tabularx}{\textwidth}{|>{\raggedright\arraybackslash}p{0.32\textwidth}|>{\raggedright\arraybackslash}X|}
\hline
\textbf{Параметр} & \textbf{Содержание} \\ \hline
Объект обследования & ООО <<\_\_\_\_>>, тематическое кафе-кофейня в стилистике Хогвартса \\ \hline
Предмет обследования & Клиентский цифровой контур: предварительное информирование, меню, события, личные мероприятия, заявка, обратная связь и цифровое продолжение тематического опыта \\ \hline
Граница обследования & Не рассматриваются касса, склад, кухня, бухгалтерия и полный контур ресторанного учета; анализ ограничен клиентскими точками контакта и данными обращений \\ \hline
Методы фиксации фактов & Аудит доступных каналов, контрольный клиентский сценарий до заявки, анализ состава информации, описание организационных ролей и AS-IS паспорта процесса личного мероприятия \\ \hline
Причина выбора показателей & У объекта отсутствует многолетняя статистика цифровых заявок, поэтому используются показатели наличия, доступности и структурированности, которые можно проверить в момент обследования \\ \hline
Использование результатов & Значения чек-листа применяются во второй главе для количественной диагностики AS-IS состояния и для обоснования проектных решений третьей главы \\ \hline
\end{tabularx}
\normalsize
\end{table}

\begin{table}[H]
\centering
\footnotesize
\caption{Чек-лист обследования каналов и клиентских данных}
\label{tab:app_channel_audit}
\setlength{\tabcolsep}{3pt}
\begin{tabularx}{\textwidth}{|>{\raggedright\arraybackslash}p{0.08\textwidth}|>{\raggedright\arraybackslash}p{0.28\textwidth}|>{\raggedright\arraybackslash}p{0.22\textwidth}|>{\centering\arraybackslash}p{0.10\textwidth}|>{\raggedright\arraybackslash}X|}
\hline
\textbf{ID} & \textbf{Проверяемый признак} & \textbf{Способ проверки} & \textbf{AS-IS} & \textbf{Пояснение к фиксации} \\ \hline
K1 & Наличие собственных управляемых цифровых каналов & Аудит каналов & 0 & Собственный сайт, Telegram-канал и карточка на картах не используются как постоянные управляемые источники \\ \hline
K2 & Доступность полного меню до визита в собственном канале & Аудит информации & 0 из 1 & Меню существует в бумажном буклете, но не доступно клиенту в собственном цифровом канале до физического визита \\ \hline
K3 & Доступность сведений о событиях и личных мероприятиях до визита & Аудит информации & 0 из 1 & Постоянный цифровой раздел с условиями событий и личных мероприятий отсутствует \\ \hline
K4 & Наличие обнаружимого цифрового сценария заявки & Контрольный клиентский сценарий & 0 из 1 & Клиент не видит единого цифрового маршрута от интереса к оформлению первичного обращения \\ \hline
K5 & Наличие единой формы заявки & Обследование процесса & 0 из 1 & Первичное обращение не нормировано в виде единой формы с заранее заданным составом полей \\ \hline
K6 & Структурированность заявки & Сопоставление фактической фиксации с минимальным набором полей & 0\% & В единой форме не фиксируются имя, телефон, дата или период, тип мероприятия, количество гостей и комментарий \\ \hline
K7 & Наличие регистра заявок со статусами & Обследование процесса & 0 из 1 & Нет единого цифрового регистра, в котором обращение получает статус и историю обработки \\ \hline
K8 & Наличие системной обратной связи & Обследование процесса & 0 из 1 & Устные реакции гостей не образуют цифровую историю отзывов и не связываются с мероприятием или клиентским действием \\ \hline
K9 & Наличие опубликованного цифрового геймифицированного профиля & Аудит клиентского канала & 0 из 1 & Тематический опыт кафе физически выражен в помещении, но не продолжен в опубликованном цифровом профиле клиента \\ \hline
K10 & Наличие правового согласия в цифровой форме заявки & Аудит формы и сценария передачи данных & 0 из 1 & Цифровой формы с персональными данными пока нет, следовательно, не оформлены цель обработки и согласие субъекта \\ \hline
\end{tabularx}
\normalsize
\end{table}

По результатам чек-листа базовое AS-IS состояние характеризуется не частичной автоматизацией, а отсутствием собственного управляемого клиентского цифрового контура. Поэтому нулевые значения показателей во второй главе не являются экспертной оценкой <<на глаз>>: они следуют из бинарной проверки наличия канала, формы, регистра, обратной связи и правового контура. Такой способ фиксации соответствует ранней стадии объекта и позволяет использовать результаты обследования как доказательную основу для проектной главы.
